% Source printed pp. 320-324 / PDF pp. 328-332.
\section{全局商叠与消解性质}\label{src:7.2}
\begin{epigraph}
你有没有停下来想一想, 然后忘了重新开始?
\epigraphauthor{小熊维尼}
\end{epigraph}
形如 $[U/G]$ 的商叠构成一类重要的代数叠. 它们的几何可通过 $U$ 的 $G$-等变几何方便地研究 (见表~\bookref{1.6.15}). 代数叠何时是商叠, 是一个有趣的问题, 它既与消解性质等几何性质有关, 也与 Brauer 映射的满射性等算术性质有关. 在 Deligne--Mumford 叠的局部结构定理 (\bookref{5.3.1}) 中, 我们证明了每个 Deligne--Mumford 叠都有保持稳定子的平展邻域, 它是有限群作用下的商叠 $[\Spec A/G]$; 稍后将对代数叠上线性约化稳定子的点证明类似结果 (见定理~\bookref{7.7.1}).
\subsection{全局商叠}\label{src:7.2.1-subsection}
\begin{definition}\label{src:7.2.1}
若存在同构 $\mathcal X\cong[U/\mathrm{GL}_n]$, 其中 $U$ 是代数空间, 则称代数叠 $\mathcal X$ 为\newterm{全局商叠}{global quotient stack}.
\end{definition}
换言之, $\mathcal X$ 是全局商叠, 当且仅当存在来自代数空间的主 $\mathrm{GL}_n$-丛 $U\to\mathcal X$, 等价地, 存在可表示态射 $\mathcal X\to B\mathrm{GL}_n$.
\begin{exercise}[重要]\label{src:7.2.2}
证明 Noether 代数叠 $\mathcal X$ 是全局商叠, 当且仅当存在 $\mathcal X$ 上的向量丛 $E$, 使对每个像闭的几何点 $x:\Spec\Bbbk\to\mathcal X$, 稳定子 $G_x$ 在纤维 $E\otimes\Bbbk$ 上忠实作用.

\textit{提示:} 使用习题~\bookref{B.1.58} 中主 $\mathrm{GL}_n$-丛与向量丛的对应.
\end{exercise}
\begin{exercise}\label{src:7.2.3}
证明 Noether 概形上一般稳定子平凡的光滑 Deligne--Mumford 叠 $\mathcal X$ 是全局商叠.

\textit{提示:} 考虑切丛和高阶 jet 丛.
\end{exercise}
\begin{exercise}\label{src:7.2.4}
若 $\mathcal X$ 是具有仿射稳定子的约化 Noether 代数叠, 证明存在分层 $\mathcal X=\coprod_i\mathcal X_i$, 其中每个 $\mathcal X_i\hookrightarrow\mathcal X$ 是局部闭子叠且是商叠. 若 $\mathcal X$ 正规, 证明 $\mathcal X_i$ 可写成商叠 $[U_i/\mathrm{GL}_n]$, 其中 $U_i$ 是拟射影概形.

\textit{提示:} 构造凝聚层, 使每个稳定子在纤维上忠实作用, 并考虑它是向量丛的开区域. 最后一个结论应用 Sumihiro 线性化定理 (\bookref{B.1.33}).
\end{exercise}
在习题~\bookref{5.3.2} 中, 我们已看到 Deligne--Mumford 叠可写成商叠 $[U/G]$, 其中 $G$ 为有限群, $U$ 为仿射概形 (分别地, 概形, 代数空间), 当且仅当它有来自仿射概形 (分别地, 概形, 代数空间) 的有限平展覆盖. 类似地, 有:
\begin{exercise}\label{src:7.2.5}
设 $\mathcal X\to\mathcal Y$ 是 Noether 代数叠之间满射, 平坦且射影的态射. 若 $\mathcal X$ 是商叠, 证明 $\mathcal Y$ 是商叠.
\end{exercise}
\begin{remark}[拟射影叠]\label{src:7.2.6}
设 $\mathcal X$ 是域 $\Bbbk$ 上分离有限型 Deligne--Mumford 叠, 粗模空间为 $\pi:\mathcal X\to X$. Kresch 的定理 \cite[定理 5.3]{Kre09} 表明, $X$ 拟射影且 $\mathcal X$ 是全局商叠, 当且仅当 $\mathcal X\cong[U/\mathrm{GL}_n]$, 其中 $U$ 拟射影. 而且, 若 $\Bbbk$ 的特征为 $0$, 则它还等价于 (a) $X$ 拟射影且 $\mathcal X$ 有生成凝聚层 $\mathcal V$, 即对 $\mathcal X$ 上每个拟凝聚层 $F$, $\pi_*\pi^*\mathcal H\!om_{\mathcal O_{\mathcal X}}(\mathcal V,F)\otimes_{\mathcal O_{\mathcal X}}\mathcal V\to F$ 满射; 或等价于 (b) 存在局部闭浸入 $\mathcal X\hookrightarrow\mathcal W$, 其中 $\mathcal W$ 是 $\Bbbk$ 上具有射影粗模空间的光滑固有 Deligne--Mumford 叠. 另见 \cite{BOW24} 中的推广.
\end{remark}
\begin{remark}[非商叠]\label{src:7.2.7}
给出具有仿射稳定子但不是全局商叠的代数叠的例子并不容易. 例如, 尚不知道每个分离 Deligne--Mumford 叠是否都是全局商叠. Noether 代数叠是全局商叠的一个必要条件是它具有拟仿射对角, 因为每个拟分离代数空间都有拟仿射对角 (推论~\bookref{5.5.8}). Kresch 在 \cite[命题 5.2]{Kre13} 中证明, 参数化亏格 $0$ 且至多有 $m\geq2$ 个结点的预稳定曲线的模叠 $\mathcal M_0^{\leq m}$ 不具有拟仿射对角. 因此 $\mathcal M_0^{\leq m}$ 不是全局商叠. 另一个必要条件是 $\mathcal X$ 有足够多向量丛 (习题~\bookref{7.2.2}), 这也可用于构造反例 (见习题~\bookref{7.2.8}). 还可用 Brauer 映射不满射来构造非商叠 (见习题~\bookref{7.4.52}).
\end{remark}
\begin{exercise}\label{src:7.2.8}
对域 $\Bbbk$, 考虑 $\mathbb P^1=\mathbb P^1_{\Bbbk}$ 上秩 $r\geq2$ 的分裂环面 $T=\mathbb G_{m,\mathbb P^1}^r$. 证明存在非平凡对合 $\alpha:T_0=\mathbb G_{m,\Bbbk}^r\to\mathbb G_{m,\Bbbk}^r=T_\infty$, 将其粘合得到射影结点三次曲线 $C=\mathbb P^1/(0\sim\infty)$ 上的非分裂环面 $T'$, 使 $BT'$ 没有向量丛. 推出 $BT'$ 不是全局商叠.
\end{exercise}
\subsection{消解性质与 Totaro 定理}\label{src:7.2.2-subsection}
商叠性质与消解性质密切相关.
\begin{definition}\label{src:7.2.9}
若每个凝聚层都是某个向量丛的商, 则称 Noether 代数叠具有\newterm{消解性质}{resolution property}.
\end{definition}
域上的光滑或拟射影概形具有消解性质. 更一般地, 有``丰沛线丛族''的概形具有消解性质, 从而每个具有仿射对角的 Noether 正规 $\mathbb Q$-析因概形都有消解性质 \cite{BS03}.
\begin{proposition}\label{src:7.2.10}
设 $G$ 是域 $\Bbbk$ 上的仿射代数群, 作用于拟射影 $\Bbbk$-概形 $U$. 假设存在带 $G$-作用的丰沛线丛 $L$ (例如 $U$ 拟仿射, $L=\mathcal O_U$). 则 $[U/G]$ 具有消解性质.
\end{proposition}
\begin{proof}
线丛 $L$ 对应于 $[U/G]$ 上关于态射 $p:[U/G]\to BG$ 相对丰沛的线丛 $\mathcal L$. 对 $[U/G]$ 上的凝聚层 $F$, 自然映射
\[
\mathcal L^{-N}\otimes p^*p_*(\mathcal L^N\otimes F)\twoheadrightarrow F
\]
在 $N\gg0$ 时满射. 推前 $p_*(\mathcal L^N\otimes F)$ 是 $BG$ 上的拟凝聚层, 即 $G$-表示, 可写成有限维 $G$-表示 $V_i$ 的并 (\bookref{B.1.20}(1)). 因而得到满射 $\operatorname{colim}_i(\mathcal L^{-N}\otimes p^*V_i)\twoheadrightarrow F$. 由于 $F$ 凝聚, 当 $i\gg0$ 时 $\mathcal L^{-N}\otimes p^*V_i\twoheadrightarrow F$ 满射.
\end{proof}
\begin{exercise}\label{src:7.2.11}
若 $U$ 是 $\mathbb Z$ 上带 $\mathrm{GL}_n$-作用的拟射影概形, 且 $U$ 有带 $\mathrm{GL}_n$-作用的丰沛线丛, 证明 $[U/\mathrm{GL}_n]$ 具有表示性质. (进一步的推广见 \cite[\S2]{Tho87}.)
\end{exercise}
Totaro 定理给出逆命题.
\begin{theorem}[Totaro 定理]\label{src:7.2.12}
设 $\mathcal X$ 是正规 Noether 代数叠, 闭点处的稳定子群仿射. 则下列条件等价:
\begin{enumerate}[(1)]
\item $\mathcal X$ 具有消解性质.
\item $\mathcal X\cong[U/\mathrm{GL}_n]$, 其中 $U$ 拟仿射.
\end{enumerate}
若 $\mathcal X$ 还是域 $\Bbbk$ 上的有限型叠, 则上述条件还等价于
\begin{enumerate}[(3)]
\item $\mathcal X\cong[\Spec A/G]$, 其中 $G$ 是 $\Bbbk$ 上的光滑仿射代数群.
\end{enumerate}
特别地, $\mathcal X$ 有仿射对角.
\end{theorem}
\begin{remark}\label{src:7.2.13}
Gross 将这一定理推广到了非正规叠 \cite{Gro17}.
\end{remark}
\begin{proof}
由于 $\mathcal X$ 正规, Sumihiro 线性化定理 (\bookref{B.1.33}) 表明 $\mathcal O_U$ 有 $\mathrm{GL}_n$-作用 (分别地, $\mathcal O_{\Spec A}$ 有 $G$-作用), 从而习题~\bookref{7.2.11} (分别地, 命题~\bookref{7.2.10}) 给出 (2)$\Rightarrow$(1) (分别地, (3)$\Rightarrow$(1)).

为证明 (3)$\Rightarrow$(2), 只需找到 $\Bbbk$ 上的忠实表示 $G\hookrightarrow\mathrm{GL}_N$, 使 $\mathrm{GL}_N/G$ 拟仿射. 事实上, 此时 $[\Spec A/G]\cong[(\Spec A\times^G\mathrm{GL}_N)/\mathrm{GL}_N]$ (习题~\bookref{4.4.19}), 且 $\Spec A\times^G\mathrm{GL}_N$ 在 $\mathrm{GL}_N/G$ 上仿射. 先选取任意忠实表示 $G\subseteq\mathrm{GL}_n$. 由 \bookref{B.1.19}(7), $\mathrm{GL}_n/G$ 拟射影, 因而存在 $\mathrm{GL}_n$-表示 $V$ 及 $\Bbbk$-点 $x\in\mathbb P(V)$, 其稳定子为 $G$. 对 $\mathrm{GL}_n\times\mathbb G_m$ 在 $\mathbb A(V)$ 上的作用 (其中 $\mathbb G_m$ 通过伸缩作用), $x$ 的提升 $\widetilde x\in\mathbb A(V)$ 的稳定子为 $G$. 由 $g\mapsto g\widetilde x$ 定义的映射 $(\mathrm{GL}_n\times\mathbb G_m)/G\hookrightarrow\mathbb A(V)$ 是局部闭浸入, 因而 $(\mathrm{GL}_n\times\mathbb G_m)/G$ 拟仿射. 在自然嵌入 $\mathrm{GL}_n\times\mathbb G_m\hookrightarrow\mathrm{GL}_{n+1}$ 下, 商 $\mathrm{GL}_{n+1}/(\mathrm{GL}_n\times\mathbb G_m)$ 仿射 (有时称为``Stiefel 流形''). 复合 $BG\to B(\mathrm{GL}_n\times\mathbb G_m)\to B\mathrm{GL}_{n+1}$ 拟仿射, 因而 $\mathrm{GL}_{n+1}/G$ 及 $U:=\Spec A\times^G\mathrm{GL}_N$ 也拟仿射.

反之, 为证明 (2)$\Rightarrow$(3), 先选取 $\mathrm{GL}_n$-等变开浸入 $U\hookrightarrow\Spec A$, 其中 $\Spec A$ 是 $\Bbbk$ 上有限型仿射概形. 事实上, 态射 $p:[U/\mathrm{GL}_n]\to B\mathrm{GL}_n$ 拟仿射, 故 $[U/\mathrm{GL}_n]\to\mathcal S\!pec_{B\mathrm{GL}_n}p_*\mathcal O_{[U/\mathrm{GL}_n]}$ 是开浸入. 将 $p_*\mathcal O_{[U/\mathrm{GL}_n]}=\operatorname{colim}_\lambda\mathcal A_\lambda$ 写成有限型 $\mathcal O_{B\mathrm{GL}_n}$-代数的余极限, 极限方法 (\bookref{B.3}) 表明当 $\lambda\gg0$ 时 $[U/\mathrm{GL}_n]\to\mathcal S\!pec_{B\mathrm{GL}_n}\mathcal A_\lambda$ 是开浸入. 设 $Z\subseteq\Spec A$ 是 $U$ 的约化补, 选取 $\mathrm{GL}_n$-等变态射 $f:\Spec A\to\mathbb A^r$ 使 $f^{-1}(0)=Z$ (命题~\bookref{B.1.21}(2)). 这给出仿射态射 $U\to\mathbb A^r\setminus0$. 补 $\mathbb A^r\setminus0$ 可实现为商 $\mathrm{GL}_r/H$, 其中 $H\cong\mathbb G_a^{r-1}\rtimes\mathrm{GL}_{r-1}\subseteq\mathrm{GL}_r$ 是最后一行为 $(0,\ldots,0,1)$ 的矩阵组成的子群. 在 $\mathrm{GL}_n$-等变 Cartesian 图
\[
\begin{tikzcd}[ampersand replacement=\&]
P\arrow[r]\arrow[d]\&\mathrm{GL}_r\arrow[d]\\
U\arrow[r]\arrow[ur,phantom,"{\square}"]\&\mathbb A^r\setminus0,
\end{tikzcd}
\]
中, $P$ 在 $\mathrm{GL}_r$ 上仿射, 因而仿射. 利用等价 $[U/\mathrm{GL}_n]\cong[P/(\mathrm{GL}_n\times H)]$, 得到结论.

剩下证明 (1)$\Rightarrow$(2). 先证明存在向量丛 $E$, 使每个稳定子群 $G_x$ 在纤维 $E\otimes\kappa(x)$ 上忠实作用. 由习题~\bookref{7.2.2}, 这蕴涵 $\mathcal X=[U/\mathrm{GL}_n]$, 其中 $n=\operatorname{rk}E$, $U=\operatorname{Fr}(E)$ 为标架丛. 对每个闭点 $x\in\mathcal X$, 设 $i_x:\mathcal G_x\hookrightarrow\mathcal X$ 是剩余 gerbe 的嵌入 (命题~\bookref{4.5.16}). 选取代表 $x$ 的有限型点 $\overline x:\Spec\Bbbk\to\mathcal X$, 得到诱导的有限态射 $p:BG_{\overline x}\to\mathcal G_x$. 由于 $G_{\overline x}$ 是仿射代数群, 可选取忠实表示 $W$ (命题~\bookref{B.1.21}(2)). 利用 $\mathcal X$ 上的消解性质, 存在向量丛 $E$ 及满射 $E\twoheadrightarrow(i_x\circ p)_*W$. $E$ 的核子群 $S_E\subseteq I_{\mathcal X}$, 即惯性叠中参数化在 $E$ 上平凡作用的元素的子群叠, 在 $x$ 上平凡. 若 $F$ 是另一个向量丛, 则 $S_{E\oplus F}\subseteq S_E$ 是闭子群. 由于 $I_{\mathcal X}$ 是 Noether 的, 可用 Noether 归纳法选出核子群 $S_E$ 平凡的向量丛 $E$.

由于 $\mathcal X$ 正规, $U$ 也正规, 可应用习题~\bookref{5.6.4}, 使 $U$ 是某个有限群 $H$ 在正规概形 $U'$ 上作用的粗模空间. 设 $g:U'\to U$ 是商态射, $\{U_i'\}$ 是有限仿射覆盖, 约化补为 $Z_i'=U'\setminus U_i'$. 则 $F:=g_*(\bigoplus_iI_{Z_i'})$ 是 $U$ 上的凝聚层, 且 $g^*F\twoheadrightarrow\bigoplus_iI_{Z_i'}$ 满射. 由于 $q:U\to\mathcal X$ 拟仿射, $q^*q_*F\twoheadrightarrow F$ 满射; 将 $q_*F$ 写成凝聚层的余极限, 可找到 $\mathcal X$ 上凝聚层 $G$ 及满射 $q^*G\to F$. 由于 $\mathcal X$ 具有消解性质, 可使 $G$ 为向量丛. 这给出满射 $g^*q^*G\twoheadrightarrow g^*F\twoheadrightarrow\bigoplus_iI_{Z_i'}$. 设 $V=\operatorname{Fr}(G)$ 为标架丛, 考虑 Cartesian 图
\[
\begin{tikzcd}[ampersand replacement=\&]
U'_V\arrow[r]\arrow[d,"{\beta}"']\&U_V\arrow[r]\arrow[d]\&V\arrow[d]\\
U'\arrow[r,"{g}"']\arrow[ur,phantom,"{\square}"]\&U\arrow[r,"{q}"']\arrow[ur,phantom,"{\square}"]\&\mathcal X,
\end{tikzcd}
\]
其中 $g$ 是主 $H$-丛, $q$ 是主 $\mathrm{GL}_n$-丛, $V\to\mathcal X$ 是主 $\mathrm{GL}_m$-丛 (其中 $m=\operatorname{rk}G$). 由于 $G$ 在 $V$ 上的拉回平凡, 拉回 $\beta^*(\bigoplus_iI_{Z_i'})$ 由整体截面生成. 这意味着 $\beta^{-1}(Z_i')$ 由 $U'_V$ 上的整体函数定义, 且补 $\beta^{-1}(U_i')$ 可由形如 $\{f\neq0\}$ 的仿射开集覆盖, 其中 $f\in\Gamma(U'_V,\mathcal O_{U'_V})$. 换言之, $\mathcal O_{U'_V}$ 丰沛, $U'_V$ 是拟仿射概形. 由于 $\beta:U'_V\to U_V$ 是有限群 $H$ 的商, $U_V$ 也拟仿射 (习题~\bookref{5.2.14}). 可写成 $\mathcal X\cong[\operatorname{Fr}(E\oplus F)/\mathrm{GL}_{n+m}]$, 其中标架丛 $\operatorname{Fr}(E\oplus F)$ 可与 $U_V\times^{\mathrm{GL}_n\times\mathrm{GL}_m}\mathrm{GL}_{n+m}$ 等同; 它拟仿射, 因为 $U_V$ 拟仿射, 商 $\mathrm{GL}_{n+m}/(\mathrm{GL}_n\times\mathrm{GL}_m)$ 仿射. 另见 \cite[定理 1.1]{Tot04}.
\end{proof}
